% GENERATED FILE. Edit canonical lessons, then run npm run generate:books.
% canonical-source-hash: fnv1a64:5ac4d678ec31ff14
% canonical-lessons: PT-C389-crime, PT-C389-criminoso, PT-C389-corrupcao, PT-C389-sancao, PT-C389-multar

\chapter{O crime não compensa}
\label{ch:pt-o-crime-nao-compensa}
\hlchaptermodality{389}

\begin{chapteropening}
\textbf{By the end of this chapter:} \emph{I can talk about crime and corruption, and the fines and sanctions that follow.}

Complete the last lesson of chapter 389: I can talk about crime and corruption, and the fines and sanctions that follow.
\end{chapteropening}

\section[o crime]{o crime — the crime}

\label{lesson:PT-C389-crime}



\begin{quote}
Before the new one: say the Portuguese for the conviction, then the Portuguese for the punishment.
\end{quote}

\subsection*{You'll want to know: o crime}
\textbf{o crime} — \textquotedblleft{}the crime\textquotedblright{}.

\begin{quote}
O crime aconteceu de madrugada.
\end{quote}

\begin{quote}
\emph{The crime happened in the early hours.}
\end{quote}

\emph{Cometer um crime} — to commit a crime.

\begin{cousinweb}
\emph{crime} was borrowed from French \emph{crime} (EtymDB), from Latin \emph{crīmen} \textquotedblleft{}an accusation, a charge, a crime\textquotedblright{} (Lewis \& Short).
\end{cousinweb}

\begin{grammarlens}[title={What you've built}]
One more word for telling, explaining and describing at B1.
\end{grammarlens}

\subsection*{Your turn}
\begin{itemize}
\raggedright
  \item \emph{Say it:} \emph{o crime}
  \item \emph{Say it:} \emph{o crime}, once more
  \item \emph{Say it:} a sentence of your own with \emph{o crime}
  \item \emph{Recall:} say the Portuguese for the conviction, then the Portuguese for the punishment, then say \emph{o crime} again
\end{itemize}

\begin{tcolorbox}[breakable,colback=teal!4,colframe=teal!35!black,title={Before you move on}]
What does o crime mean? (\textquotedblleft{}The crime\textquotedblright{}.) Say it once more.
\end{tcolorbox}
\section[o criminoso]{o criminoso — the criminal}

\label{lesson:PT-C389-criminoso}



\begin{quote}
Before the new one: say the Portuguese for the punishment, then the Portuguese for the crime.
\end{quote}

\subsection*{You'll want to know: o criminoso}
\textbf{o criminoso} — \textquotedblleft{}the criminal\textquotedblright{}.

\begin{quote}
A polícia apanhou o criminoso.
\end{quote}

\begin{quote}
\emph{The police caught the criminal.}
\end{quote}

Also an adjective: \emph{um ato criminoso}.

\begin{cousinweb}
\emph{criminoso} is a learned borrowing from Latin \emph{crīminōsus} \textquotedblleft{}reproachful, criminal\textquotedblright{} (EtymDB, Lewis \& Short), built on \emph{crīmen} \textquotedblleft{}a charge, a crime\textquotedblright{}.
\end{cousinweb}

\begin{grammarlens}[title={What you've built}]
One more word for telling, explaining and describing at B1.
\end{grammarlens}

\subsection*{Your turn}
\begin{itemize}
\raggedright
  \item \emph{Say it:} \emph{o criminoso}
  \item \emph{Say it:} \emph{o criminoso}, once more
  \item \emph{Say it:} a sentence of your own with \emph{o criminoso}
  \item \emph{Recall:} say the Portuguese for the punishment, then the Portuguese for the crime, then say \emph{o criminoso} again
\end{itemize}

\begin{tcolorbox}[breakable,colback=teal!4,colframe=teal!35!black,title={Before you move on}]
What does o criminoso mean? (\textquotedblleft{}The criminal\textquotedblright{}.) Say it once more.
\end{tcolorbox}
\section[a corrupção]{a corrupção — the corruption}

\label{lesson:PT-C389-corrupcao}



\begin{quote}
Before the new one: say the Portuguese for the crime, then the Portuguese for the criminal.
\end{quote}

\subsection*{You'll want to know: a corrupção}
\textbf{a corrupção} — \textquotedblleft{}the corruption\textquotedblright{}.

\begin{quote}
Há corrupção em muitos países.
\end{quote}

\begin{quote}
\emph{There is corruption in many countries.}
\end{quote}

\emph{Corrupto} — corrupt.

\begin{cousinweb}
\emph{corrupção} is a learned borrowing from Latin \emph{corruptiō} \textquotedblleft{}a corrupting, a spoiling\textquotedblright{} (EtymDB), from \emph{corrumpere}, built on \emph{rumpere} \textquotedblleft{}to break\textquotedblright{} (Lewis \& Short).
\end{cousinweb}

\begin{grammarlens}[title={What you've built}]
One more word for telling, explaining and describing at B1.
\end{grammarlens}

\subsection*{Your turn}
\begin{itemize}
\raggedright
  \item \emph{Say it:} \emph{a corrupção}
  \item \emph{Say it:} \emph{a corrupção}, once more
  \item \emph{Say it:} a sentence of your own with \emph{a corrupção}
  \item \emph{Recall:} say the Portuguese for the crime, then the Portuguese for the criminal, then say \emph{a corrupção} again
\end{itemize}

\begin{tcolorbox}[breakable,colback=teal!4,colframe=teal!35!black,title={Before you move on}]
What does a corrupção mean? (\textquotedblleft{}The corruption\textquotedblright{}.) Say it once more.
\end{tcolorbox}
\section[a sanção]{a sanção — the sanction, the penalty}

\label{lesson:PT-C389-sancao}



\begin{quote}
Before the new one: say the Portuguese for the criminal, then the Portuguese for the corruption.
\end{quote}

\subsection*{You'll want to know: a sanção}
\textbf{a sanção} — \textquotedblleft{}the sanction, the penalty\textquotedblright{}.

\begin{quote}
A empresa sofreu sanções por poluir o rio.
\end{quote}

\begin{quote}
\emph{The company faced sanctions for polluting the river.}
\end{quote}

\emph{Aplicar uma sanção} — to impose a sanction.

\begin{cousinweb}
\emph{sanção} is a learned borrowing from Latin \emph{sānctiō} \textquotedblleft{}a decree, a penalty clause\textquotedblright{} (EtymDB, Lewis \& Short), from \emph{sancīre} \textquotedblleft{}to make sacred, to decree\textquotedblright{}.
\end{cousinweb}

\begin{grammarlens}[title={What you've built}]
One more word for telling, explaining and describing at B1.
\end{grammarlens}

\subsection*{Your turn}
\begin{itemize}
\raggedright
  \item \emph{Say it:} \emph{a sanção}
  \item \emph{Say it:} \emph{a sanção}, once more
  \item \emph{Say it:} a sentence of your own with \emph{a sanção}
  \item \emph{Recall:} say the Portuguese for the criminal, then the Portuguese for the corruption, then say \emph{a sanção} again
\end{itemize}

\begin{tcolorbox}[breakable,colback=teal!4,colframe=teal!35!black,title={Before you move on}]
What does a sanção mean? (\textquotedblleft{}The sanction, the penalty\textquotedblright{}.) Say it once more.
\end{tcolorbox}
\section[multar]{multar — to fine}

\label{lesson:PT-C389-multar}



\begin{quote}
Before the new one: say the Portuguese for the corruption, then the Portuguese for the sanction.
\end{quote}

\subsection*{You'll want to know: multar}
\textbf{multar} — \textquotedblleft{}to fine\textquotedblright{}.

\begin{quote}
A polícia multou-me por excesso de velocidade.
\end{quote}

\begin{quote}
\emph{The police fined me for speeding.}
\end{quote}

\emph{Ser multado} — to be fined; \emph{a multa} — the fine.

\begin{cousinweb}
\emph{multar} comes from Latin \emph{multāre} \textquotedblleft{}to punish, to fine\textquotedblright{} (EtymDB), on \emph{multa} \textquotedblleft{}a fine\textquotedblright{} (Lewis \& Short).
\end{cousinweb}

\begin{grammarlens}[title={What you've built}]
One more word for telling, explaining and describing at B1.
\end{grammarlens}

\subsection*{Your turn}
\begin{itemize}
\raggedright
  \item \emph{Say it:} \emph{multar}
  \item \emph{Say it:} \emph{multar}, once more
  \item \emph{Say it:} a sentence of your own with \emph{multar}
  \item \emph{Recall:} say the Portuguese for the corruption, then the Portuguese for the sanction, then say \emph{multar} again
\end{itemize}

\begin{tcolorbox}[breakable,colback=teal!4,colframe=teal!35!black,title={Before you move on}]
What does multar mean? (\textquotedblleft{}To fine\textquotedblright{}.) Say it once more.
\end{tcolorbox}
